\section{Digital phantom}\label{sec:phantom}

\subsection{Phantom and simulated acquisition}

The phantom is a $256\times256$ axial-slice-like arrangement of tissue classes
(\cref{fig:phantom}): an outer CSF rim, cortical grey matter with a gyrated boundary, white
matter, two lateral ventricles, two deep grey-matter nuclei and a white-matter lesion, in
total \Voxels{} voxels (\cref{tab:tissues}). Each class carries 3\,T-typical relaxation values,
modulated by a smooth random texture of $3\%$ (log scale) so that tissues are not uniform.
Smooth fields model the transmit field (centre brightening, $\Bone=0.87$--$1.15$, so that about a
quarter of the voxels have $\alpha_2>360^\circ$ and a ring of voxels has $\alpha_2\approx360^\circ$), the
off-resonance (a quadratic background of $\pm20$\,Hz and a frontal hot spot of about $+50$\,Hz),
the receive sensitivity ($\pm15\%$) and the receive phase.

\begin{table}[!htbp]
\centering
\caption{Phantom tissue classes (before the $3\%$ texture).}
\label{tab:tissues}
\small
\begin{tabular}{lrrrrr}
\toprule
Tissue & $\Tone$ (ms) & $\Ttwo$ (ms) & $\Ttwos$ (ms) & PD & voxels\\
\midrule
CSF & 4000 & 1500 & 900 & 1.00 & 5614\\
GM & 1350 & 90 & 60 & 0.82 & 14\,782\\
WM & 850 & 66 & 50 & 0.69 & 16\,196\\
deep GM & 1130 & 58 & 40 & 0.75 & 1104\\
lesion & 1500 & 120 & 80 & 0.85 & 180\\
\bottomrule
\end{tabular}
\end{table}

\begin{figure}[!htbp]
\centering
\includegraphics[width=0.92\linewidth]{phantom_truth.png}
\caption{True parameter maps. The dashed contour marks $\alpha_2=330^\circ\Bone=360^\circ$.}
\label{fig:phantom}
\end{figure}

Images are generated voxel by voxel with the forward model \eqref{eq:signal} (512
isochromats; no partial volume within a voxel) and complex Gaussian noise with
$\sigma=1/\mathrm{SNR}(\Mzero)$, for SNR$(\Mzero)=1000$ and $300$ ($\Mzero$ ranges from 0.6 to 1.2). Two
acquisitions are simulated (\cref{fig:images}):
\begin{itemize}
\item \textbf{Scenario A}: both scans at full resolution, the per-voxel setting of
  \cref{sec:information,sec:ml,sec:magnitude};
\item \textbf{Scenario B}: the published design, in which scan 2 covers only the central
  $25\%$ of the phase-encoding lines ($k_y$), with zero filling. (In 3D the published design
  keeps $25\%\times25\%$ of $k_y$--$k_z$, a factor 16; in 2D only one phase-encoding axis can be
  truncated, a factor 4.) Noise is added before truncation, i.e.\ the same per-sample noise
  with fewer samples.
\end{itemize}

\begin{figure}[!htbp]
\centering
\includegraphics[width=\linewidth]{phantom_images.png}
\caption{Simulated MPME images (magnitude, SNR$(\Mzero)=1000$). The scan-2 images show the dark
ring where $\alpha_2^{\rm eff}=360^\circ$ and the signal vanishes; bottom right: the low-resolution scan 2 of
Scenario B, with truncation ringing.}
\label{fig:images}
\end{figure}

\subsection{Reconstructions}

All fits use the exact Jacobian in single precision with 128 isochromats, 25 iterations,
the $\Bone\le540/330$ bound and four starts (\cref{alg:ml}), processed in chunks of 4096 voxels.
\begin{itemize}
\item \textbf{Scenario A.} The analytic inverse (voxel-wise flip angle, no smoothing);
  complex ML; magnitude least squares with the FID-sign constraint. For the last two, a
  fifth \emph{spatial-restart} start is added after the first pass: the first-pass $\Bone$ map is
  smoothed by a robust polynomial fit, and each voxel is restarted from
  $(\Mzero\kappa,\ B^{+}_{1,\rm smooth},\ \Tone\kappa^2)$ with $\kappa=\Bone/B^{+}_{1,\rm smooth}$, i.e.\ moved along the
  near-degenerate direction of \cref{prop:scaling}; the lowest cost is kept. The neighbourhood
  only proposes a start, so every estimate remains a per-voxel likelihood maximum.
\item \textbf{Scenario B.} The paper pipeline (\cref{alg:twostage}): the flip angle from the
  two low-resolution scans by Eq.~15, a degree-4 robust polynomial fit, then $\Tone$, $\Ttwo$ and
  $\Mzero$ from full-resolution scan 1 by Eqs.~1, 16 and 17. The ML pipeline: $\Bone$ from the two
  low-resolution scans by complex ML, the same polynomial fit, then complex ML on
  full-resolution scan 1 with $\Bone$ fixed.
\end{itemize}

\paragraph{Why the spatial restart.} Without it, 324 voxels (SNR 1000) had $|\log(\hat\Tone/\Tone)|>0.3$;
321 of them lie within $3^\circ$ of $\alpha_2=360^\circ$, where scan 2 carries almost no signal and the problem
approaches the one-scan degeneracy of \cref{cor:onescan}. In 222 of the 324 the final cost
exceeded the cost at the true parameters (optimiser failures: all four starts derive from
estimators that are unreliable at $360^\circ$). The smoothed $\Bone$ map is accurate to $0.03\%$ there,
and the restart removes all optimiser failures (no voxel ends above the true-parameter cost).
The remaining large errors (53 at SNR 1000, 50 of them within $3^\circ$ of $360^\circ$) all have a cost at
or below the true one: they are genuine maxima of the likelihood, a property of the protocol
at $\alpha_2\approx360^\circ$ rather than of the estimator.

\subsection{Results}

\begin{table}[!htbp]
\centering
\caption{Phantom, SNR$(\Mzero)=1000$: median and SD of $\log(\hat\theta/\theta)$ over each tissue (bias and spread), and the fraction of voxels with an undefined estimate. Truth includes a $3\%$ smooth texture.}
\label{tab:roi1000}
\footnotesize
\setlength{\tabcolsep}{3.5pt}
\begin{tabular}{llcccccc}
\toprule
Tissue & Method & $\Tone$ median & $\Tone$ SD & $\Ttwo$ median & $\Ttwo$ SD & $\Bone$ median & undef.\,\%\\
\midrule
WM & A, analytic inverse & +0.006 & 0.298 & -0.000 & 0.145 & -0.000 & 0.0\\
 & A, complex ML & +0.001 & 0.094 & +0.001 & 0.041 & +0.000 & 0.0\\
 & A, magnitude LS + FID & -0.000 & 0.124 & +0.003 & 0.041 & +0.000 & 0.0\\
 & B, paper pipeline & +0.001 & 0.101 & -0.000 & 0.145 & +0.001 & 0.0\\
 & B, ML pipeline & +0.001 & 0.073 & -0.000 & 0.055 & -0.000 & 0.0\\
\midrule
GM & A, analytic inverse & +0.003 & 0.403 & +0.001 & 0.160 & +0.000 & 0.0\\
 & A, complex ML & +0.001 & 0.042 & +0.001 & 0.030 & +0.000 & 0.0\\
 & A, magnitude LS + FID & +0.002 & 0.042 & +0.000 & 0.030 & +0.000 & 0.0\\
 & B, paper pipeline & +0.043 & 0.117 & +0.001 & 0.160 & -0.015 & 0.0\\
 & B, ML pipeline & +0.001 & 0.065 & +0.001 & 0.051 & +0.000 & 0.0\\
\midrule
deep GM & A, analytic inverse & +0.011 & 0.321 & -0.009 & 0.141 & -0.000 & 0.0\\
 & A, complex ML & -0.000 & 0.115 & +0.000 & 0.041 & +0.000 & 0.0\\
 & A, magnitude LS + FID & -0.002 & 0.131 & +0.003 & 0.041 & +0.000 & 0.0\\
 & B, paper pipeline & +0.006 & 0.154 & -0.009 & 0.141 & -0.001 & 0.0\\
 & B, ML pipeline & +0.001 & 0.084 & -0.000 & 0.052 & -0.000 & 0.0\\
\midrule
lesion & A, analytic inverse & -0.004 & 0.119 & -0.009 & 0.167 & -0.000 & 0.0\\
 & A, complex ML & -0.007 & 0.050 & -0.005 & 0.031 & -0.000 & 0.0\\
 & A, magnitude LS + FID & -0.008 & 0.050 & -0.004 & 0.031 & -0.000 & 0.0\\
 & B, paper pipeline & -0.008 & 0.118 & -0.009 & 0.167 & +0.001 & 0.0\\
 & B, ML pipeline & +0.005 & 0.059 & -0.004 & 0.050 & +0.000 & 0.0\\
\midrule
CSF & A, analytic inverse & +0.015 & 0.650 & -0.019 & 0.802 & -0.000 & 0.1\\
 & A, complex ML & -0.001 & 0.079 & -0.000 & 0.080 & +0.000 & 0.0\\
 & A, magnitude LS + FID & -0.001 & 0.079 & -0.000 & 0.080 & +0.000 & 0.0\\
 & B, paper pipeline & +0.085 & 0.356 & -0.019 & 0.802 & -0.015 & 0.1\\
 & B, ML pipeline & -0.004 & 0.099 & +0.003 & 0.114 & -0.000 & 0.0\\
\bottomrule
\end{tabular}
\end{table}

\begin{table}[!htbp]
\centering
\caption{Phantom, SNR$(\Mzero)=300$: median and SD of $\log(\hat\theta/\theta)$ over each tissue (bias and spread), and the fraction of voxels with an undefined estimate. Truth includes a $3\%$ smooth texture.}
\label{tab:roi300}
\footnotesize
\setlength{\tabcolsep}{3.5pt}
\begin{tabular}{llcccccc}
\toprule
Tissue & Method & $\Tone$ median & $\Tone$ SD & $\Ttwo$ median & $\Ttwo$ SD & $\Bone$ median & undef.\,\%\\
\midrule
WM & A, analytic inverse & +0.015 & 0.652 & +0.003 & 0.771 & -0.002 & 1.3\\
 & A, complex ML & +0.000 & 0.230 & +0.005 & 0.138 & +0.000 & 0.0\\
 & A, magnitude LS + FID & -0.009 & 0.250 & +0.026 & 0.139 & +0.000 & 0.0\\
 & B, paper pipeline & -0.017 & 0.452 & +0.003 & 0.771 & +0.001 & 1.3\\
 & B, ML pipeline & -0.006 & 0.225 & +0.008 & 0.176 & -0.000 & 0.0\\
\midrule
GM & A, analytic inverse & +0.015 & 0.913 & +0.004 & 0.879 & +0.000 & 1.6\\
 & A, complex ML & +0.002 & 0.139 & +0.004 & 0.102 & +0.000 & 0.0\\
 & A, magnitude LS + FID & +0.013 & 0.137 & +0.001 & 0.100 & +0.002 & 0.0\\
 & B, paper pipeline & +0.027 & 0.508 & +0.004 & 0.879 & -0.016 & 1.6\\
 & B, ML pipeline & -0.000 & 0.208 & +0.009 & 0.168 & +0.000 & 0.0\\
\midrule
deep GM & A, analytic inverse & -0.018 & 0.728 & -0.023 & 0.762 & -0.002 & 5.4\\
 & A, complex ML & +0.002 & 0.287 & +0.003 & 0.170 & +0.000 & 0.0\\
 & A, magnitude LS + FID & -0.026 & 0.301 & +0.033 & 0.143 & +0.000 & 0.0\\
 & B, paper pipeline & -0.048 & 0.656 & -0.023 & 0.762 & -0.000 & 5.4\\
 & B, ML pipeline & -0.003 & 0.276 & +0.004 & 0.173 & +0.000 & 0.0\\
\midrule
lesion & A, analytic inverse & -0.032 & 0.747 & -0.028 & 0.933 & -0.001 & 1.7\\
 & A, complex ML & -0.026 & 0.161 & -0.006 & 0.100 & -0.000 & 0.0\\
 & A, magnitude LS + FID & -0.030 & 0.160 & +0.007 & 0.100 & -0.000 & 0.0\\
 & B, paper pipeline & -0.037 & 0.593 & -0.028 & 0.933 & +0.001 & 1.7\\
 & B, ML pipeline & -0.009 & 0.179 & -0.014 & 0.160 & -0.000 & 0.0\\
\midrule
CSF & A, analytic inverse & +0.047 & 1.011 & -0.054 & 1.228 & -0.010 & 17.8\\
 & A, complex ML & -0.015 & 0.275 & +0.003 & 0.280 & +0.000 & 0.0\\
 & A, magnitude LS + FID & -0.015 & 0.275 & +0.002 & 0.280 & +0.000 & 0.0\\
 & B, paper pipeline & +0.050 & 0.757 & -0.054 & 1.228 & -0.020 & 17.8\\
 & B, ML pipeline & -0.035 & 0.343 & +0.031 & 0.393 & +0.000 & 0.0\\
\bottomrule
\end{tabular}
\end{table}



\begin{figure}[!htbp]
\centering
\includegraphics[width=\linewidth]{phantomA_T1_snr1000.png}
\caption{Scenario A, SNR$(\Mzero)=1000$: $\Tone$ maps and log-ratio error maps.}
\label{fig:phA_T1}
\end{figure}

\begin{figure}[!htbp]
\centering
\includegraphics[width=\linewidth]{phantomA_B1_snr1000.png}
\caption{Scenario A, SNR$(\Mzero)=1000$: $\Bone$ maps and log-ratio error maps.}
\label{fig:phA_B1}
\end{figure}

\begin{figure}[!htbp]
\centering
\includegraphics[width=\linewidth]{phantomA_T2_snr1000.png}
\caption{Scenario A, SNR$(\Mzero)=1000$: $\Ttwo$ maps and log-ratio error maps.}
\label{fig:phA_T2}
\end{figure}

\paragraph{Scenario A.} The maps reproduce the voxel-wise analysis
(\cref{fig:phA_T1,fig:phA_B1,fig:phA_T2}, \cref{tab:roi1000}). The analytic inverse is
unbiased in the median but noisy: at SNR 1000 its $\Tone$ spread is $0.30$ in white matter and
$0.40$ in grey matter, against $0.042$ for complex ML in grey matter, and its errors grow
towards the periphery (lower $\Bone$) and form a sharp ring at $\alpha_2=360^\circ$. Complex ML and
magnitude least squares with the FID sign are nearly indistinguishable, with negligible bias
in every tissue; their white-matter spread ($0.094$ and $0.124$) is inflated by the genuine
outliers on the $360^\circ$ ring, which lies in white matter, while the grey-matter spread ($0.042$)
is close to the voxel-wise bound. No voxel farther than $3^\circ$ from $360^\circ$ ends on the wrong branch
for any method.

\begin{figure}[!htbp]
\centering
\includegraphics[width=\linewidth]{phantomB_snr1000.png}
\caption{Scenario B (scan 2 at $25\%$ of $k_y$), SNR$(\Mzero)=1000$. Top: true $\Bone$; low-resolution
$\Bone$ from Eq.~15; its polynomial fit; low-resolution $\Bone$ from complex ML. Bottom: $\Tone$ and its error
for the two pipelines.}
\label{fig:phB}
\end{figure}

\paragraph{Scenario B.} The low-resolution flip-angle map of the paper pipeline shows
truncation ringing, errors at the $360^\circ$ ring and outliers at the periphery
(\cref{fig:phB}); the polynomial fit absorbs much of this but remains biased by $-1.5\%$ in grey
matter (up to $15\%$ at the edge). Through the error transfer $\partial\log\Tone/\partial\log\Bone\approx-2$ of the
scan-1 $\Tone$ step (\cref{tab:cond}) this produces a $\Tone$ bias of $+4.3\%$ in grey matter and
$+8.5\%$ in CSF; with the true $\Bone$, the same $\Tone$ step is unbiased ($-0.06\%$ in grey matter). The
ML pipeline's smoothed map is accurate to $0.2\%$ everywhere and its $\Tone$ is unbiased in every
tissue. Its $\Tone$ spread is larger than in Scenario A in grey matter ($0.065$ vs $0.042$), as
expected from the scan-1-only bound ($57.1$ vs $40.5$), but smaller in white matter ($0.073$ vs
$0.094$), because a smooth, fixed $\Bone$ removes the $360^\circ$ ambiguity. Its $\Ttwo$ spread is larger than
in Scenario A because stage 2 uses only scan 1.

\paragraph{SNR 300.} \Cref{tab:roi300} gives the same statistics at SNR$(\Mzero)=300$
(image SNR about 26). The analytic inverse breaks down: its $\Tone$ spread is
$0.65$ in white matter and $0.91$ in grey matter, and $\Tone$ is undefined in $1.3$--$1.7\%$ of tissue
voxels and $17.8\%$ of CSF voxels. Complex ML and magnitude least squares with the FID sign
remain nearly unbiased, with spreads of $0.14$ (grey matter) to $0.28$ (CSF), no undefined
estimates and no wrong-branch voxels away from the $360^\circ$ ring (\cref{fig:phA_T1_300}). In
Scenario B the paper pipeline has spreads of $0.45$--$0.51$ in tissue and the same undefined
fractions, whereas the ML pipeline is nearly unbiased with spreads of $0.21$--$0.23$; its
grey-matter spread is $1.5$ times that of Scenario A, close to the bound ratio $57.1/40.5=1.4$.
At this SNR, errors above $0.3$ in log units are expected from noise alone in long-$\Tone$ tissue
(the local bound for CSF is of that order), and only 4 of \Voxels{} complex-ML voxels end above
the true-parameter cost.

\begin{figure}[!htbp]
\centering
\includegraphics[width=\linewidth]{phantomA_T1_snr300.png}
\caption{Scenario A, SNR$(\Mzero)=300$: $\Tone$ maps and log-ratio error maps.}
\label{fig:phA_T1_300}
\end{figure}

\subsection{Uncertainty maps}

\begin{figure}[!htbp]
\centering
\includegraphics[width=\linewidth]{phantom_uncertainty_snr1000.png}
\caption{Complex ML, Scenario A, SNR$(\Mzero)=1000$: predicted SD of $\log\Tone$ (local CRLB at the
estimate), absolute error, and the distribution of the standardised error $z$ against $N(0,1)$
($z$ clipped to $\pm5$ for display).}
\label{fig:phunc}
\end{figure}

The local CRLB evaluated at the estimate gives a predicted standard deviation for every voxel
at the cost of one Jacobian (\cref{fig:phunc}). The map shows the structure expected from
\cref{sec:information}: higher uncertainty in CSF and long-$\Tone$ tissue, in deep grey matter, at
low $\Bone$, and on the $360^\circ$ ring. Standardised errors are close to $N(0,1)$. For complex ML in
Scenario A, $95\%$ intervals cover the truth in \CalAcovTThousand\% of voxels at SNR 1000
(\CalAcovTThreeHundred\% at SNR 300); the SD of $z$ is \CalAzTThousand{} (\CalAzTThreeHundred),
above 1 only because of the genuine outliers at $360^\circ$, which a local bound cannot describe.
For the ML pipeline of Scenario B (with $\Bone$ treated as known), coverage is
\CalBcovTThousand\% and \CalBcovTThreeHundred\% and the SD of $z$ is \CalBzTThousand{} and
\CalBzTThreeHundred{}: slightly conservative at SNR 300, consistent with the constraint
$\Rtwop\ge0$ being active in many voxels at that SNR (a constrained estimator can have a variance
below the interior bound, \cref{sec:ml}).

\subsection{Run times}

\begin{table}[!htbp]
\centering
\caption{Reconstruction times for the phantom (voxels in the mask; laptop CPU, 8 threads, implicit Jacobian, float64, 128 isochromats, 40 iterations).}
\label{tab:phantomtime}
\small
\begin{tabular}{lrr}
\toprule
Step & SNR 1000 & SNR 300\\
\midrule
A: analytic inverse & 3.0 s & 3.2 s\\
init (model fit) & 8.4 s & 9.1 s\\
A: complex ML (4 starts) & 131.0 s & 139.8 s\\
A: complex ML spatial restart & 34.1 s & 34.2 s\\
A: magnitude LS + FID sign (4 starts) & 120.9 s & 124.8 s\\
A: magnitude LS spatial restart & 30.2 s & 29.7 s\\
A: uncertainty maps & 2.6 s & 2.6 s\\
B: paper pipeline & 3.9 s & 4.2 s\\
B: ML pipeline & 194.8 s & 212.1 s\\
\bottomrule
\end{tabular}
\end{table}


A full set of reconstructions for one SNR level takes about 9 minutes on the laptop CPU
(\cref{tab:phantomtime}); each full-slice fit with four starts takes about 2 minutes, and the
uncertainty maps a few seconds.
